\documentclass[11pt]{article}
\usepackage[margin=1in]{geometry}
\usepackage{amsmath,amssymb,amsfonts,bm}
\usepackage{booktabs,multirow,array}
\usepackage{graphicx}
\usepackage{xcolor}
\usepackage{hyperref}
\usepackage{cleveref}
\usepackage{siunitx}
\usepackage{enumitem}
\usepackage{microtype}
\usepackage{authblk}
\emergencystretch=3em
\hypersetup{colorlinks=true,linkcolor=blue!60!black,citecolor=blue!60!black,urlcolor=blue!60!black}
\newcommand{\E}{\mathbf{e}}
\newcommand{\norm}[1]{\left\lVert #1 \right\rVert}
\title{\textbf{CrossPiezo: Databases Agree on Promising Regions but Cannot Resolve the Elite Tail in Piezoelectric Screening}}
\author[1]{Anonymous Author(s)}
\affil[1]{Anonymous Institution}
\date{30 July 2026}
\begin{document}
\maketitle
\begin{abstract}
High-throughput piezoelectric databases are widely used to train and benchmark machine-learning models.
Here we reframe cross-database comparison as a question of \emph{screening resolution}: do the Materials Project and JARVIS databases agree on broad promising regions of chemical space, yet fail to reproducibly resolve the elite tail of candidates?
Using the frozen CrossPiezo-Invariant-v1 benchmark of 573 strict structure-matched pairs (P0) and 207 tighter matches (P2), we compute coordinate-invariant response scalars F1, F3 and F4 and sweep the screened quantile from 1\% to 50\%.
The area under the chance-adjusted concordance curve (AUC-Concordance) is 5.15\% for F1 on P0; the lower confidence bound of adjusted overlap first exceeds zero only at 13\%.
Quantile-normalization does not restore tail stability, and simpler scalar controls such as unit-cell volume and band gap are markedly more cross-source consistent than the piezoelectric response.
Source-robust portfolios outperform single-source selection on worst-source recall and NDCG, but this is two-source robustness, not physical validation.
We release hash-bound tables and reproducibility instructions.
\end{abstract}
\section{Introduction}
Piezoelectric response tensors are central to electromechanical materials design, yet costly to compute with density-functional perturbation theory.
The Materials Project (MP) and JARVIS projects have released thousands of computed tensors \cite{dejong2015,choudhary2020}, and recent equivariant neural networks report strong in-source prediction accuracy \cite{yan2024,dong2025}.
Improved agreement with a held-out subset of one database does not guarantee that a candidate is robustly ranked across independently generated high-throughput datasets \cite{hegde2023}.
\section{Screening-resolution framing}
We ask whether two databases can identify broad high-response regions while failing to resolve the elite tail.
For each panel, metric and quantile $q=1\%\ldots50\%$, we report observed and chance-adjusted top-$q$ Jaccard overlap, exact hypergeometric null, bootstrap confidence intervals, and source rank displacement.
\section{Results}
The screening-resolution curves show low adjusted overlap at small $q$ and a slow rise toward larger $q$.
AUC-Concordance and $q_\text{consensus}$ are reported in \Cref{tab:auc_summary}.
Quantile normalization leaves Kendall $\tau$ nearly unchanged, indicating that the instability is not merely a monotonic scale shift.
Source-robust portfolios improve worst-source recall and NDCG relative to JARVIS-only or MP-only selection.
\section{Discussion}
The weak elite-tail resolution is a statement about database definitions and processed fields, not about the physical correctness of either source.
A third computational protocol or experimental adjudication set would be required to identify physical truth.
\section{Data availability}
Hash-bound result tables, panel membership and reproducibility instructions are released in the CrossPiezo repository.
\bibliographystyle{unsrt}
\bibliography{CrossPiezo_ScreeningResolution_references}
\end{document}
